\documentclass[10pt]{article}

\usepackage[margin=0.75in]{geometry}
\usepackage{amsmath,amsthm,amssymb}
\usepackage{xcolor}
\usepackage{cancel}
\usepackage{graphicx}
\usepackage{changepage}
\usepackage{circuitikz}
\usepackage{pgfplots}
\usepackage{physics}
\usepackage{hyperref}
\usepackage{siunitx}
\usepackage{fontspec}
\usepackage{relsize}
\usepackage{subfig}
\usepackage{minted}
\usepackage{pdfpages}
\usepackage{todonotes}
\usepackage{multicol, multirow, booktabs}
\usepackage[breakable]{tcolorbox}
\usepackage[inline]{enumitem}

\theoremstyle{definition}
\newtheorem{problem}{Problem}
\newtheorem{soln}{Solution}

\pgfplotsset{compat=newest}
\usetikzlibrary{lindenmayersystems}
\usetikzlibrary{arrows}
\usetikzlibrary{calc}
\usetikzlibrary{positioning, fit}
\usetikzlibrary{3d, perspective}

\definecolor{incolor}{HTML}{303F9F}
\definecolor{outcolor}{HTML}{D84315}
\definecolor{cellborder}{HTML}{CFCFCF}
\definecolor{cellbackground}{HTML}{F7F7F7}
\newcommand{\ui}{\hat{i}}
\newcommand{\uj}{\hat{j}}
\newcommand{\uk}{\hat{k}}
\newcommand{\ux}{\hat{\mathbf{x}}}
\newcommand{\uy}{\hat{\mathbf{y}}}
\newcommand{\uz}{\hat{\mathbf{z}}}
\newcommand{\uv}[1]{\hat{\bm{{#1}}}}
\newcommand{\pr}[1]{{#1^\prime}}
\newcommand{\pd}[2][]{{\frac{\partial^{#1}}{\partial {{#2}^{#1}}}}}
\newcommand{\dif}[2][]{{\frac{d^{#1}}{d{{#2}^{#1}}}}}
\pgfdeclarelayer{background}  
\pgfsetlayers{background,main}
\AtBeginDocument{\RenewCommandCopy\qty\SI}
\newcommand{\justif}[2]{&{#1}&\text{#2}}
\DeclareMathOperator\Arg{Arg}
\DeclareMathOperator{\sgn}{sgn}

\makeatletter
\newcommand{\boxspacing}{\kern\kvtcb@left@rule\kern\kvtcb@boxsep}
\makeatother
\newcommand{\prompt}[4]{
    \ttfamily\llap{{\color{#2}[#3]:\hspace{3pt}#4}}\vspace{-\baselineskip}
}

\newcommand{\thevenin}[2]{
  \begin{center}
    \begin{circuitikz} \draw
      (0,0) -- (2,0) to[battery1, l_=$V_{Th}\eq#1$] (2,2) 
      to[resistor, l_=$R_{Th}\eq#2$] (0,2)
      ;
      \draw [o-] (-.07,2.079);
      \draw [o-] (-.07,0.079);
    \end{circuitikz}
  \end{center}
}

\newcommand{\norton}[2]{
  \begin{center}
    \begin{circuitikz} \draw
      (0,0) -- (3,0) to[american current source, l_=$I_{N}\eq#1$] (3,2) -- (0,2) (2,0)
      to[resistor, l=$R_{N}\eq#2$] (2,2)
      ;
      \draw [o-] (-.07,2.079);
      \draw [o-] (-.07,0.079);
    \end{circuitikz}
  \end{center}
}

\newcommand{\highlight}[1]{{\colorbox{yellow}{#1}}}

\newcommand{\ti}[1]{\widetilde{#1}}

\newfontface{\Kaufmann}{Kaufmann}
\DeclareTextFontCommand{\kf}{\Kaufmann}
\newcommand{\scriptr}{\fontsize{12pt}{12pt}\kf{r}}

\newfontface{\KaufmannB}{Kaufmann Bd BT}
\DeclareTextFontCommand{\kfb}{\KaufmannB}
\newcommand{\bscriptr}{\fontsize{12pt}{12pt}\kfb{r}}

\newcommand{\bv}[1]{\mathbf{#1}}

\title{Math 3160H: Assignment II}
\author{Jeremy Favro (0805980) \\ Trent University, Peterborough, ON, Canada}
\date{\today}

\begin{document}
\maketitle

% PROBLEM 1
\begin{problem}
A simple pendulum initially has the length $\ell_0$ and makes an angle $\theta_0$ with the vertical line. It is
then released from this position. If the length of the pendulum increases with time $t$ according to
$\ell=\ell_0+\epsilon t$, where $\epsilon$ is a small constant, then the angle $\theta$ that the pendulum makes with the vertical
line, assuming that the oscillations are small, is the solution of the initial value problem
$$\left(\ell_0+\epsilon t\right)\frac{d^2\theta}{dt^2}+2\epsilon\frac{d\theta}{dt}+g\theta=0,\quad \theta(0)=\theta_0,\quad \frac{d\theta}{dt}(0)=0$$
where $g$ is the acceleration due to gravity. Find the solution of the above initial value problem.
\end{problem}
\begin{soln}
Consider the change of variables 
$$t=\ell=\ell_0+\epsilon t$$. This gives us 
$d\ell=\epsilon dt$
and so 
$$\frac{d\theta}{dt}=\epsilon\frac{d\theta}{d\ell}$$ and 
$$\epsilon^2\frac{d^2\theta}{d\ell^2}$$
so the ODE becomes 
$$\ell \epsilon^2\frac{d^2\theta}{d\ell^2}+2\epsilon^2\frac{d\theta}{d\ell}+g\theta=0$$
which we can divide through by $\epsilon^2$ to see that 
$$\ell\frac{d^2\theta}{d\ell^2}+2\frac{d\theta}{d\ell}+\frac{g}{\epsilon^2}\theta=0$$
which is very nearly the Bessel equation, we just need to fiddle a bit more. We can multiply by $\ell$ to get
$$\ell^2\frac{d^2\theta}{d\ell^2}+2\ell\frac{d\theta}{d\ell}+\frac{g\ell}{\epsilon^2}\theta=0$$
\end{soln}
\end{document}